

\chapter{\parbox{0.84\textwidth}{Логічні оператори (and, or, not)}}

Прості умовні вирази (операції відношення) можна об'єднувати у складні \textbf{складені умови} за допомогою \textbf{логічних операторів} (булевих зв'язок). В Python використовуються три базові логічні оператори: \textbf{\texttt{and}, \texttt{or} та \texttt{not}}.

\vspace{0.5em}

\noindent\textbf{1. Кон'юнкція (\texttt{and} — логічне «І»)}

Оператор \textbf{\texttt{and}} повертає \textbf{\texttt{True}} тоді й лише тоді, коли \textbf{обидва} операнди є істинними. Якщо хоча б один із виразів хибний (\textbf{\texttt{False}}), увесь складений вираз набуває значення \textbf{\texttt{False}}.

Пріоритет оператора \textbf{ \texttt{and}} нижчий за пріоритет операторів порівняння (\texttt{>}, \texttt{<}, \texttt{==} тощо), тому у виразах на зразок:
\begin{py}{Перевірка діапазону}
if length > 8 and length < 20:
    print("Довжина відповідає вимогам")
\end{py}
\noindent дужки навколо окремих порівнянь писати не обов'язково.

\vspace{0.5em}

\noindent\textbf{2. Диз'юнкція (\texttt{or} — логічне «АБО»)}

Оператор \textbf{\texttt{or}} повертає \textbf{\texttt{True}}, якщо \textbf{хоча б один} із операндів є істинним (або обидва одночасно). Вираз набуває значення \textbf{\texttt{False}} лише тоді, коли всі складові частини є хибними.

\begin{py}{Приклад диз'юнкції}
if age < 18 or age >= 60:
    print("Пільгова категорія")
\end{py}

\vspace{0.5em}

\noindent\textbf{3. Заперечення (\texttt{not} — логічне «НЕ»)}

Унарний оператор \textbf{\texttt{not}} інвертує (змінює на протилежне) значення булевого виразу: \textbf{\texttt{True}} перетворюється на \textbf{\texttt{False}}, а \textbf{\texttt{False}} — на \textbf{\texttt{True}}.

\begin{py}{Приклад заперечення}
if not is_valid:
    print("Помилка введення даних")
\end{py}

\vspace{0.5em}

\begin{table}[htbp]
	\centering
	\caption{Таблиця істинності для логічних операцій}
	\begin{tabular}{|c|c|c|c|c|}
		\hline
		\multicolumn{2}{|c|}{Значення операндів} & \multicolumn{3}{c|}{Результат операції} \\ \hline
		X & Y & not X & X and Y & X or Y \\ \hline
		False & False & True  & False   & False  \\ \hline
		False & True  & True  & False   & True   \\ \hline
		True  & False & False & False   & True   \\ \hline
		True  & True  & False & True    & True   \\ \hline
	\end{tabular}
\end{table}

\vspace{0.5em}

\textbf{Особливості Python:}
\begin{itemize}
	\item \textbf{Пріоритет виконання:} Найвищий пріоритет має \texttt{not}, далі йде \texttt{and}, і найнижчий — \texttt{or}. Для зміни порядку обчислення використовують круглі дужки \texttt{()}.
	\item \textbf{Ліниві обчислення (Short-circuit evaluation):} Python обчислює логічні вирази зліва направо і зупиняється, як тільки результат стає однозначним (наприклад, якщо у виразі з \texttt{and} перший операнд \texttt{False}, другий операнд навіть не обчислюється).
\end{itemize}
\chapter{\parbox{0.75\textwidth}{Генерація випадкових чисел: модуль random}}

При розв'язанні багатьох математичних задач, моделюванні ймовірнісних процесів (метод Монте-Карло) або створенні ігор виникає потреба у використанні **випадкових чисел**. У Python для цього призначений вбудований модуль \texttt{random}.

Щоб отримати доступ до функцій модуля, його необхідно імпортувати на початку програми:
\begin{py}{Імпорт модуля}
import random
\end{py}

\section*{Основні функції модуля random}

\noindent\textbf{1. Дійсні числа: \texttt{random.random()} та \texttt{random.uniform()}}

\begin{itemize}
	\item \texttt{random.random()} — повертає випадкове дійсне число $x$ у напіввідкритому інтервалі $[0.0, 1.0)$.
	\item \texttt{random.uniform(a, b)} — повертає випадкове дійсне число $x$ у відрізку $[a, b]$.
\end{itemize}

\begin{py}{Генерація дійсних чисел}
import random
# Випадкове число від 0.0 до 1.0
p = random.random()
print("Імовірність",p)

# Випадкова температура в межах [-5.0, 15.0]
temp = random.uniform(-5.0, 15.0)
print("Температура:", temp)
\end{py}

\vspace{0.5em}

\noindent\textbf{2. Цілі числа: \texttt{random.randint()} та \texttt{random.randrange()}}

\begin{itemize}
	\item \texttt{random.randint(a, b)} — повертає випадкове ціле число $n \in [a, b]$ (включаючи обидва кінці $a$ та $b$).
	\item \texttt{random.randrange(start, stop, step)} — повертає випадковий елемент із послідовності \texttt{range(start, stop, step)}.
\end{itemize}

\begin{py}{Генерація цілих чисел}
import random
# Моделювання кидка грального кубика (від 1 до 6)
dice = random.randint(1, 6)
print("Випало число:", dice)

# Випадкове парне число від 10 до 98
even_num = random.randrange(10, 100, 2)
print("Парне число:",even_num)
\end{py}

\vspace{0.5em}


\section*{Відтворюваність результатів: \texttt{random.seed()}}

Оскільки Python генерує **псевдовипадкові числа** за допомогою алгоритму (Вихр Мерсенна), послідовність обчислень залежить від початкового значення (значення ініціалізації — \textit{seed}).

Функція \texttt{random.seed(a)} фіксує початковий стан генератора. Це є критично важливим у наукових та математичних обчисленнях для **відтворюваності результатів експерименту**.

\begin{py}{Фіксація генератора}
import random

random.seed(42)  # Задаємо фіксоване зерно
print(random.randint(1, 100))  # Завжди повертатиме один і той самий результат при повторних запусках
\end{py}


